% Repository-only archive. This file is intentionally not included by main.tex.
% It preserves the former historical m-free optimization subsection verbatim.

\subsection{The historical \texorpdfstring{$m$}{m}-free model}\label{sec:mstar-model}

\Cref{ch:certify} reduces the directed multigraph problem at every $m$ to one $m$-free optimisation, $M^{*}(n)$, by dividing multiplicities by $m-1$. This subsection writes out the model that computed it, the theorem it settled and the record of what was run. \Cref{cor:mstar-integral} above has since replaced it, so nothing below is load-bearing.

By max-flow/min-cut, $\mathrm{maxflow}(s,t)\le1$ exactly when some $s$--$t$ cut has capacity at most one. The optimisation therefore chooses one cut per ordered pair. A binary label $x^{st}_u$ records its side, and $p^{st}_{uv}$ counts $w(u,v)$ only when $u\to v$ crosses from source to sink (\Cref{fig:cut}). This describes the exact feasible region. The archived wrapper retained an \texttt{OPTIMAL} status and primal witness but no independently replayable matching bound, so the runs are finite solver checks rather than formal certificates.

With scaled weights $w$, binary side labels $x$, and non-negative crossing weights $p$, the model is
\begin{align*}
    \text{maximise} \quad & \sum_{u \ne v} w(u, v) \\
    \text{subject to} \quad & x^{st}_s = 1, \quad x^{st}_t = 0 && \text{for all } (s,t), \\
    & p^{st}_{uv} \ge w(u, v) + x^{st}_u - x^{st}_v - 1 && \text{for all } (s,t),\ (u,v), \\
    & \sum_{u \ne v} p^{st}_{uv} \le 1 && \text{for all } (s,t), \\
    & w(s,t) + \sum_{x} \min\bigl(w(s,x), w(x,t)\bigr) \le 1 && \text{for all } (s,t), \\
    & d(v_0) \le d(v_1) \le \dots \le d(v_{n-1}), && \text{with } d(v) = \textstyle\sum_{u \ne v} \bigl(w(u,v) + w(v,u)\bigr), \\
    & w \in [0,1], \quad x \in \{0,1\}.
\end{align*}

The first line is what makes the labelling an $s$--$t$ cut rather than an arbitrary bipartition, and it is not optional. Without it the all-zero labelling $x \equiv 0$ is admissible, no ordered pair $(u,v)$ ever crosses, every $p^{st}_{uv}$ rests at $0$, and the cut constraint below is satisfied by every weight matrix whatsoever. Pinning the source side of $s$ and the sink side of $t$ removes that escape and leaves the remaining $n-2$ labels free, which is precisely the choice of an $s$--$t$ cut. The implementation supplies these two values as constants rather than as variables (\code{program/erdos915\_unified.py}, in \code{\_cut\_counting\_model}), which is the same restriction imposed one step earlier.

Because $w\le1$, the second constraint forces $p^{st}_{uv}\ge w(u,v)$ only in the crossing case $(x^{st}_u,x^{st}_v)=(1,0)$. In the other three cases its right-hand side is non-positive. The third constraint then caps the chosen cut at one. The remaining two lines are redundant accelerators: the first counts the direct arc and all two-step detours, and the second removes equivalent vertex relabellings.

The $\min$ in the two-step line is written as a minimum for readability only. A linear program cannot take one directly. It is linearised in the usual way, by an auxiliary variable $z^{sxt} \in [0,1]$ held below both arguments, $z^{sxt} \le w(s,x)$ and $z^{sxt} \le w(x,t)$, together with a binary selector $b^{sxt}$ and the two big-$M$ inequalities $z^{sxt} \ge w(s,x) - (1 - b^{sxt})$ and $z^{sxt} \ge w(x,t) - b^{sxt}$, which force $z^{sxt}$ up to whichever argument is the smaller. The displayed $\min$ is then the value of $z^{sxt}$.

\begin{figure}[ht]
\centering
\begin{tikzpicture}[>=Stealth, every node/.style={font=\small},
    plab/.style={edgelabel}]
    % the dashed cut line
    \draw[dashed, thick, color=KULblauw3b] (0,-2.3) -- (0,2.3);
    \node[color=KULblauw3b] at (0,2.65) {cut};
    \node at (-2.4,2.65) {source side ($x=1$)};
    \node at ( 2.4,2.65) {sink side ($x=0$)};
    % source-side vertices
    \node[vertex] (s)  at (-2.6, 1.1) {$s$};
    \node[vertex] (a)  at (-1.3,-1.1) {$a$};
    % sink-side vertices
    \node[vertex] (t)  at ( 2.6,-1.1) {$t$};
    \node[vertex] (b)  at ( 1.3, 1.1) {$b$};
    % non-crossing arcs (p = 0): labels pushed to the outer side, white-filled
    \draw[->, edge, color=gray] (s) -- node[plab, left=2pt]{$p=0$} (a);
    \draw[->, edge, color=gray] (b) -- node[plab, right=2pt]{$p=0$} (t);
    % crossing arcs source -> sink (p = w), highlighted
    \draw[->, edge, very thick, color=KULblauw1] (s) -- node[plab, above]{$p=w$} (b);
    \draw[->, edge, very thick, color=KULblauw1] (a) -- node[plab, below]{$p=w$} (t);
    % a back-crossing arc sink -> source (p = 0), not counted; label on the apex
    \draw[->, edge, color=gray] (b) to[bend left=28] node[plab, pos=0.30, below=1pt]{$p=0$} (a);
\end{tikzpicture}
\caption{One chosen $s$--$t$ cut. Only the two thick source-to-sink arcs contribute crossing weight $p=w$, while same-side and backward arcs have $p=0$.}
\label{fig:cut}
\end{figure}
\begin{theorem}[Directed multigraph value, small $n$]\label{thm:dir-multi-small}
For $n \le 6$, $\;M^{*}(n) = 2(n-1)$, and the optimum is attained by a $\{0,1\}$ weight matrix. Hence for all $m \ge 2$, $\;L_m^{\mathrm{dir}}(n) = 2(n - 1)(m - 1)$, attained by the double star.
\end{theorem}

Historical runs gave $M^{*}(n)=2(n-1)$ with status \texttt{OPTIMAL} for $3\le n\le6$, 
and $n=7$ did not close. Their archived output lacks a separately checkable bound or gap, 
and no theorem now depends on it: \Cref{app:proofs} proves the value for all $n,m$ by hand 
and \Cref{sec:transcripts} retains the solver record.
